% !TEX program = lualatex
% Generated by new_homework.sh — Math 118 Homework 3
\documentclass[11pt]{article}

\usepackage{fontspec}
\usepackage{microtype}
\usepackage{mathtools,amssymb,amsfonts,amsthm,bm}
\usepackage{enumitem}
\usepackage{booktabs,array,xcolor}
\usepackage{geometry,fancyhdr,setspace}
\usepackage{pdfpages}
\usepackage[hidelinks]{hyperref}
\usepackage[nameinlink,noabbrev]{cleveref}

% ---------- Page and paragraph rhythm (matches lecture/solutions templates) ----------
\geometry{margin=1in,headheight=15pt}
\setstretch{1.12}
\setlength{\parindent}{0pt}
\setlength{\parskip}{0.55em plus 0.12em minus 0.08em}
\setlength{\abovedisplayskip}{0.8em plus 0.2em minus 0.1em}
\setlength{\belowdisplayskip}{0.8em plus 0.2em minus 0.1em}
\allowdisplaybreaks[3]
\setlist{leftmargin=*,topsep=0.35em,itemsep=0.15em,parsep=0pt}

% ---------- Metadata ----------
\newcommand{\studentname}{Keshav Ramamurthy}
\newcommand{\coursename}{Math 118}
\newcommand{\courselongname}{Fourier Analysis}
\newcommand{\assignment}{Homework 3}
\newcommand{\term}{Fall 2026}

% ---------- Core notation (same as ~/.config/nvim/textemplate.tex) ----------
\newcommand{\N}{\mathbb{N}}
\newcommand{\Z}{\mathbb{Z}}
\newcommand{\Q}{\mathbb{Q}}
\newcommand{\R}{\mathbb{R}}
\newcommand{\C}{\mathbb{C}}
\newcommand{\F}{\mathbb{F}}
\newcommand{\K}{\mathbb{K}}
\newcommand{\E}{\mathbb{E}}
\newcommand{\Prob}{\mathbb{P}}
\newcommand{\1}{\mathbf{1}}
\newcommand{\eps}{\varepsilon}
\newcommand{\dd}{\,\mathrm{d}}
\newcommand{\abs}[1]{\left\lvert#1\right\rvert}
\newcommand{\norm}[1]{\left\lVert#1\right\rVert}
\newcommand{\ip}[2]{\left\langle#1,#2\right\rangle}
\newcommand{\set}[1]{\left\{#1\right\}}
\newcommand{\ceil}[1]{\left\lceil#1\right\rceil}
\newcommand{\floor}[1]{\left\lfloor#1\right\rfloor}
\newcommand{\given}{\,\middle\vert\,}
\newcommand{\indep}{\perp\!\!\!\perp}
\newcommand{\wh}[1]{\widehat{#1}}
\newcommand{\deriv}[2]{\frac{\mathrm{d}#1}{\mathrm{d}#2}}
\newcommand{\pderiv}[2]{\frac{\partial#1}{\partial#2}}
\newcommand{\lto}{\longrightarrow}
\newcommand{\st}{\text{ such that }}
\DeclareMathOperator{\Aut}{Aut}
\DeclareMathOperator{\End}{End}
\DeclareMathOperator{\Gal}{Gal}
\DeclareMathOperator{\Hom}{Hom}
\DeclareMathOperator{\Ker}{ker}
\DeclareMathOperator{\im}{im}
\DeclareMathOperator{\rank}{rank}
\DeclareMathOperator{\nullity}{nullity}
\DeclareMathOperator{\Span}{span}
\DeclareMathOperator{\tr}{tr}
\DeclareMathOperator{\diag}{diag}
\DeclareMathOperator{\spec}{spec}
\DeclareMathOperator{\ord}{ord}
\DeclareMathOperator{\supp}{supp}
\DeclareMathOperator{\diam}{diam}
\DeclareMathOperator{\Var}{Var}
\DeclareMathOperator{\Cov}{Cov}
\DeclareMathOperator*{\argmin}{arg\,min}
\DeclareMathOperator*{\argmax}{arg\,max}

% ---------- Statement environments ----------
\newtheorem{problem}{Problem}
\newtheorem{theorem}{Theorem}
\newtheorem{lemma}{Lemma}
\newtheorem{proposition}{Proposition}
\newtheorem{corollary}{Corollary}
\theoremstyle{definition}
\newtheorem{definition}{Definition}
\newtheorem{example}{Example}
\newtheorem{remark}{Remark}

% ---------- Solution machinery (matches solutions.tex / textemplate.tex) ----------
\newenvironment{solution}{\begin{proof}[Solution]}{\end{proof}}
\newenvironment{answer}{\begin{proof}[Answer]}{\end{proof}}
\renewcommand{\qedsymbol}{\(\square\)}
\newcommand{\exercise}[1]{%
  \par\goodbreak\medskip\noindent\textbf{Exercise #1}\par\nopagebreak\smallskip\nopagebreak}
\newcommand{\exref}[1]{\textsf{\textbf{Exercise~#1}}}
\newenvironment{claim}{\par\smallskip\noindent\textit{Claim.}\ }{\par\smallskip}
\newenvironment{idea}{\par\smallskip\noindent\textit{Idea.}\ }{\par\smallskip}
\newenvironment{proofsketch}{\par\smallskip\noindent\textit{Proof sketch.}\ }{\par\smallskip}

% ---------- Header ----------
\pagestyle{fancy}
\fancyhf{}
\fancyhead[L]{\coursename}
\fancyhead[C]{\assignment\ — my solutions}
\fancyhead[R]{\studentname}
\fancyfoot[C]{\thepage}

\begin{document}

% ---------- Title ----------
\begin{center}
  {\Large\sffamily\bfseries \coursename\ -- \courselongname}\par\vspace{3pt}
  {\sffamily \assignment\ — my solutions}\par\vspace{3pt}
  {\small\sffamily \term\quad\textbullet\quad \studentname}
\end{center}
\medskip
\hrule
\bigskip

\includepdf[pages=-,pagecommand={\thispagestyle{empty}}]{hw03.pdf}
\clearpage

\section*{Solutions}
\subsubsection*{Problem 1}
\begin{solution}
We are given a few data points: \[x = {0, 1, 3, 4}\]
with corresponding outputs \(y = {0, 8, 8, 20}\)
In order to continue we consider an arbitrary polynomial
\[p(x) = ax^2 + bx + c\]
To proceed, let's calculate basic values of p that will help later:
\[p(0) = c, p(1) = a + b + c, p(3) = 9a + 3b + c, p(4) = 16a + 4b + c\]
Now, we have suppose the output vector for these inputs is
\[f = \begin{pmatrix}
    0\\
     8\\
     8\\
     20
\end{pmatrix}\]
And unlike a linear regression, our but now using a \(4 \times 3 \) matrix representing the
the \(x^2, x, 1\) terms.
\[
Z =
\begin{pmatrix}
0^2 & 0 & 1 \\
1^2 & 1 & 1 \\
3^2 & 3 & 1 \\
4^2 & 4 & 1
\end{pmatrix},
\qquad
Y =
\begin{pmatrix}
0 \\
8 \\
8 \\
20
\end{pmatrix}.
\]

\[
Z^T Y
=
ZZ^T
\begin{pmatrix}
a \\
b \\
c
\end{pmatrix}.
\]
Calculating \(Z^T Y\) to start yields 
\[
Z^T Y =
\begin{pmatrix}
0^2 & 1^2 & 3^2 & 4^2 \\
0 & 1 & 3 & 4 \\
1 & 1 & 1 & 1
\end{pmatrix}
\begin{pmatrix}
0 \\
8 \\
8 \\
20
\end{pmatrix}
=
\begin{pmatrix}
0^2(0) + 1^2(8) + 3^2(8) + 4^2(20) \\
0(0) + 1(8) + 3(8) + 4(20) \\
0 + 8 + 8 + 20
\end{pmatrix}
=
\begin{pmatrix}
400 \\
112 \\
36
\end{pmatrix}.
\]
% And notice that \(\langle f - p, 1\rangle\)
Now, let's calculate \(ZZ^T\) which is: 
\[
Z^T Z =
\begin{pmatrix}
0^2 & 1^2 & 3^2 & 4^2 \\
0 & 1 & 3 & 4 \\
1 & 1 & 1 & 1
\end{pmatrix}
\begin{pmatrix}
0^2 & 0 & 1 \\
1^2 & 1 & 1 \\
3^2 & 3 & 1 \\
4^2 & 4 & 1
\end{pmatrix}
\]

\[
=
\begin{pmatrix}
0^4 + 1^4 + 3^4 + 4^4
&
0^3 + 1^3 + 3^3 + 4^3
&
0^2 + 1^2 + 3^2 + 4^2
\\
0^3 + 1^3 + 3^3 + 4^3
&
0^2 + 1^2 + 3^2 + 4^2
&
0 + 1 + 3 + 4
\\
0^2 + 1^2 + 3^2 + 4^2
&
0 + 1 + 3 + 4
&
1 + 1 + 1 + 1
\end{pmatrix}
=
\begin{pmatrix}
338 & 92 & 26 \\
92 & 26 & 8 \\
26 & 8 & 4
\end{pmatrix}.
\]
Now we can solve for constants \(a, b\) and \(c\).
We have :
\[
\begin{pmatrix}
338 & 92 & 26\\
92 & 26 & 8\\
26 & 8 & 4
\end{pmatrix}
\begin{pmatrix}
a\\
b\\
c
\end{pmatrix}
=
\begin{pmatrix}
400\\
112\\
36
\end{pmatrix}.
\]
This yields the constants \(a = \frac{2}{3}, b = \frac{4}{3}, c = 2\)
so our best fit quadratic is \[\frac{2}{3}x^2 + \frac{4}{3}x + 2\]
% Write your proof, calculation, or argument here.
\end{solution}
\subsubsection*{Problem 2}
\begin{solution}
    In order to show that \(f_n \rightarrow 0\) in \(L^2[0, 1]\), 
    Let's consider \(f(k) = 0\), Then notice this proving this function goes to 0 is
    analogous to showing for some arbitrary \(\epsilon > 0\), \[\exists N  \text{ such
    that } n > N \implies  ||f_n||_{L^2}= \sqrt{\langle f_n, f_n\rangle} < \epsilon \]
    So this means we should calculate the inner product of \(f_n\) on itself, notice that
    is 
    \[\langle f_n, f_n \rangle  = \int_{0}^{\frac{1}{n}}(1^2)dt + \int_{\frac{1}{n}}^{1}
    0^2 dt = \frac{1}{n}\]
    So \(||f_n|| = \sqrt{\frac{1}{n}}\)
    \par So when we consider the limit as \(n \rightarrow \infty\), this is pretty trivial
    to show and can formulate it as follows:
    \par    For an  arbitrary \(\epsilon > 0\), we have \[\exists N = \left\lceil\frac{1}{e^2}\right\rceil
    \text{ such that } n > N \implies ||f_n|| < \epsilon \]
    so we are done.
    Let's prove that it does not uniformly converge, and this is pretty direct.
\par    For a given \(\epsilon > 0\), notice that we'd need to have \[||f_n|| < \epsilon, \forall t \in
[0,1] \]
But notice pretty straightforwardly  if \(t \in [0,\frac{1}{n})\) that \(f_n(t) = 1\), and
thus: \(|f_n(t) - 0| \not< \epsilon\) for any \(0 < \epsilon < 1\)so it can't be uniformly convergent
% Write your proof, calculation, or argument here.
\end{solution}
\subsubsection*{Problem 3}
\begin{solution}
In order to do number 3, we can do the same as above. In order to determine \(L^2\)
convergence, consider an arbitrary \(\epsilon > 0\), and we'd need to have 
\[
\exists N \in \mathbb{Z}^{+} \text{ such that } n \ge N \implies ||f_n - f||_{L^2} < \epsilon.
\]
In this case, we are again examining convergence to 0, so let's consider
\[
||f_n - 0||_{L^2}
= \sqrt{\int_{0}^{1} \langle f_n(t),f_n(t)\rangle\,dt}.
\]
On \(0 \leq t \leq \frac{1}{n^2}\), we have
\[
\langle f_n(t),f_n(t)\rangle
= \langle \sqrt{n},\sqrt{n}\rangle
= n,
\]
so
\[
||f_n - 0||_{L^2}
= \sqrt{\int_{0}^{\frac{1}{n^2}} n\,dt
+ \int_{\frac{1}{n^2}}^{1} 0\,dt}.
\]
Evaluating the RHS integral gives
\[
\int_{0}^{\frac{1}{n^2}} n\,dt
= \left.nt\right|_{0}^{1/n^2}
= \frac{1}{n}.
\]
Thus,
\[
||f_n - 0||_{L^2}
= \sqrt{\frac{1}{n}}
= \frac{1}{\sqrt{n}}.
\]
Notice we can really easily bound this limit, and that means just setting
\[
N = \left\lceil \frac{1}{\epsilon^2}\right\rceil + 1,
\]
we have that
\[
n >= N \implies ||f_n - 0||_{L^2} < \epsilon
\]
and thus it converges to 0 in \(L^2\).

Now to show that \(f_n(0)\) does not converge uniformly to zero, notice that
\[
f_n(0) = \sqrt{n}.
\]
Thus \(f_n(0)\) (for non zero n) does not converge to 0, and \(f_n(0) \to \infty\).
% Write your proof, calculation, or argument here.
\end{solution}
\subsubsection*{Problem 4}
\begin{solution}
Let's try to recall normal uniform converge and then \(L^2\) convergence on \([a, b]\) and then expand it to an infinite
interval. 
We are given a function \(f_n\) that converges uniformly to \(f\) in this case we can write
it, given an arbitrary \(\epsilon > 0\).
\[\exists N \text{ such that } n \ge N \implies  ||f_n(t) - f|| < \epsilon \text{ } \forall t \in
\mathbb{R}^{+}\]
In contrast, the statement for \(L^2\) convergence, can take the limit of the \(L^2\)
product on \([0, k]\) for \(k \in \mathbb{R}^{+}\) and then take the limit as \(k \rightarrow
\infty\).
This is
\[
\lim_{k \to \infty} \int_{0}^{k} {||f_n(t) - f||_{L^2}}^2\,dt.
\]
However, notice that at any point \(t \in \mathbb{R}^{+}\), we have that \[||f_n(t) - f|| <
\epsilon \implies ||f_n(t) - f||^2 < \epsilon^2 \] 
So we can bound this integral \[\lim_{k\to\infty} \int_{0}^{k} {||f_n(t) -
f||_{L^2}}^2\, dt \le \lim_{k\to\infty} \int_{0}^{k} \epsilon^2 = \lim_{k\to\infty} k
\epsilon^2\]
Since \(\lim_{k\to\infty} k\epsilon^2 = +\infty\) we are motivated to find a
counterexample.
We'll use a similar piecewise function in order to create something such that uniformly
it converges to 0, but make it such that it has an \(L^2\) inner product that doesn't
converge to 0. To do so we consider \[
f_n(t)=
\begin{cases}
\frac{1}{n}, & 0 \leq t \leq n^2,\\
0, & t > n^2.
\end{cases}
\]
In this function, uniform convergence is just for arbitrary \(\epsilon > 0\), consider 
\(N \ge \lceil\frac{1}{\epsilon}\rceil + 1\), and we'd then have that \[n \ge N \implies ||f_n(t) - 0|| <\epsilon \] so it uniformly converges to 0.
\par This is really convenient for us as when considering the \(L^2\) inner product,
notice that \[\int_{0}^{n^2} \left(\frac{1}{n}\right)^2 dt  + \int_{n^2}^{\infty} 0 dt  = \left.\frac{1}{n^2} \right\rvert_{0}^{n^2} = 1 \ne 0\] 
    so we have an effective counterexample.
% Write your proof, calculation, or argument here

\end{solution}
\subsubsection*{Problem 5}

\begin{solution}
    Let's just do this by showing \((1) \implies (2) \implies (3) \implies (1)\)
    \par In order to do so, let's consider the column vectors \(a_i\) and we are given
    \(\|a_i\| = 1, \langle a_i, a_j \rangle = 0\)
    by definition.
    \par How does this show \(AA^* = I\)
    \par
    \noindent\textit{Claim.} \((A^*A)_{ii} = 1\) for \(i \in \{1, 2,\cdots, n\}\)
    \begin{proof}
        Notice \((AA^*)_{ii}\) is a lot of computation, but we notice that if we can show \(A^*\) is a left inverse of \(A\), since they are square matrices, we'd have
        \[
        A^*A = I, \qquad A^*A = I \implies AA^* = I.
        \]
        And this is super nice because \((A^*A)_{ij}\) is just \(\langle a_i, a_j\rangle\),
        and this is good because for an identity matrix, we just have to show
        \[
        (A^*A)_{ii} = 1 \text{ and for } i \ne j, \text{ that } (A^*A)_{ij} = 0.
        \]
        But notice that \((A^*A)_{ii} = \langle a_i, a_i \rangle = 1\) due to our
        vectors being normalized, and
        \[
        (A^*A)_{ij} = \langle a_i, a_j\rangle = 0,
        \]
        so we've shown \(A^*A = I\), and we have shown \((1) \implies (2).\)

        \par Now we'll show that \((2) \implies (3)\).
        This is pretty direct from the adjoint formula, namely:
        \[
        \|Ax\|^2 = \langle Ax, Ax\rangle
        = \langle x, A^*Ax\rangle
        = \langle x, x\rangle
        = \|x\|^2.
        \]
        Thus, \(\|Ax\| = \|x\|\).

        \par Now let's show that \((3) \implies (1)\).
        \par In order to do so, we write \(A = (a_1, a_2, \cdots, a_n)\) for these column vectors,
        and similarly now denote \(x = (x_1, x_2, \cdots, x_n)\) for scalars \(x_i \in \mathbb{C}\).

        \par Notice here we have that:
        \[
        \|x\|^2 = \sum_{i=1}^{n} |x_i|^2.
        \]
        If we consider the value of the expression
        \[
        \left\lVert Ax\right\rVert
        = \left\lVert \sum_{i=1}^{n} a_i x_i\right\rVert.
        \]
        This sucks, so let's try phrasing it with the adjoint:
        \[
        \langle Ax, Ax\rangle = \langle x, x\rangle
        \implies \langle x, A^*A x\rangle = \langle x, x\rangle
        \implies \langle x, (A^*A-I)x\rangle = 0
        \]
        for all \(x \in \mathbb{C}^{n}\), or \(A^*A = I\), which is \((2)\). How does this show \((1)?\)

        \par We've shown so far \((1) \implies (2) \iff (3)\), and wish to show \((2) \implies (1)\).

        \par In order for \(A^*A = I\), we can use the same exact idea as showing \((1) \implies (2)\).

        \par This is true only if \(\langle a_i, a_i \rangle = 1\) and \(\langle a_i, a_j \rangle = 0
        \text{ if } i \ne j.\)

        \par We claim that these conditions are sufficient for the \(a_i\) being an orthonormal basis
        in \(\mathbb{C}^n\).

        \par Let's first show that these are linearly independent.
        This would be true if \(\sum_{i=1}^{n} c_i a_i = 0 \implies c_i = 0\) for constants \(c_i\).
        We take the inner product of this with one arbitrary \(a_k\), to get
        \[
        \left\langle \sum_{i=1}^{n} c_i a_i, a_k \right\rangle
        = \langle 0, a_k \rangle
        \iff \sum_{i=1}^{n} c_i \langle a_i, a_k \rangle = 0.
        \]
        This is really convenient as we've established when \(i \ne k\), \(\langle a_i, a_k \rangle
        = 0\), so this equation reduces to
        \[
        c_k \langle a_k, a_k \rangle = c_k \cdot 1 = 0
        \implies c_k = 0.
        \]

        \par This inner product can be taken for each \(k \in \{1, 2, \cdots, n\}\), which then
        shows that each \(c_i\) must be \(0\), establishing linear independence.
        Along with this, the columns are already orthonormal as
        \(\|a_i\|^2 = \langle a_i, a_i\rangle = 1\), so we are done(and it spans $C^n$
        since there are n linearly independent vectors) and we've shown $$(1) \iff (2)
        \iff (3)$$
    \end{proof}
\end{solution}
\subsubsection*{Problem 6}
\begin{solution}
    \par We are given the Fourier sine series
\[
f(x)
=
\sum_{k=1}^{\infty}
\left(
\frac{2}{\pi}
\int_0^\pi \cos(x)\sin(kx)\,dx
\right)\sin(kx).
\]
Notice we can manipulate the integrand using
\[
\cos(x)\sin(kx)
=
\frac{1}{2}\left(\sin((k+1)x)+\sin((k-1)x)\right).
\]
Plugging this back into the sum, we get
\[
f(x)
=
\sum_{k=1}^{\infty}
\left(
\frac{1}{\pi}
\int_0^\pi
\left(\sin((k+1)x)+\sin((k-1)x)\right)\,dx
\right)\sin(kx).
\]
Integrating, we get
\[
f(x)
=
\sum_{k=1}^{\infty}
\frac{1}{\pi}
\left[
-\frac{\cos((k+1)x)}{k+1}
-\frac{\cos((k-1)x)}{k-1}
\right]_{0}^{\pi}
\sin(kx).
\]
Evaluating this, for $k \ne 1$, this is equivalent to $$b_k =
\frac{1}{\pi}\left(1 - (-1)^{k + 1}\right) \left(\frac{1}{k+1}+\frac{1}{k-1}\right)$$
And for the case $k = 1$, we have that $b_1 = 0$ because 
$$b_1 = \frac{2}{\pi}\int_{0}^{\pi }\cos x \sin x = \frac{1}{\pi}\int_{0}^{\pi}
\sin(2x) = \frac{1}{\pi} \left.\left(-\frac{1}{2}\cos(2x)\right)\right|_{0}^{\pi} = 0$$
    For now we'll assume $k \ne 1$ and when is $k$ is odd, $k = 2k' + 1$, notice $(1 - (-1)^{2k' + 1 + 1}) = 0$
    so $b_k = b_{2k' + 1} = 0$
    Then, when $k$ is even where $k = 2k'$, we have that this evaluates to $$b_{2k'} =
    b_{k} =
    \frac{4k}{\pi (k+1)(k-1)}$$
    So our final fourier series is 
    $$f(x) = \cos(x) = \sum_{k=1}^{\infty}b_k \sin(kx) = \sum_{k' = 1}^{\infty}
    \frac{8k'}{\pi (2k' + 1)(2k'- 1)}  \sin(2k'x) \text{ for } 0 < x < \pi$$ (where this is reindexed cause
    we are only have non zero even indices)
\end{solution}
\end{document}
